\section[L-参数]{L-参数（2 月 8 日）}

\srcpage{138}

\subsection{设定}

设 $E$ 是非阿基米德局部域，$G/E$ 是约化群。

局部 Langlands 对应预测存在一个典范映射
\[
\{\text{不可约光滑 }G(E)\text{-表示}\}\longrightarrow\{\text{L-参数}\},
\]
我们把它记作 $\pi\mapsto\phi_\pi$。

通常人们用 $\CC$-系数来建立它。注意存在一个「典范的」$\sqrt q\in\CC$，它的选取隐含在局部 Langlands 对应之中——在一般的代数闭域上，必须作此选取才能得到对应。

\subsection{L-群}

对 $G/E$，我们有上次讨论过的对偶群 $\hat G/\ZZ$。它带有 $\Gal(\ol E/E)$ 的作用，该作用经过一个有限商 $Q$ 分解。

\begin{definition}
Langlands $L$-群是 $\prescript{L}{}{G}:=\hat G\rtimes Q$，视为 $\ZZ$ 上的代数群。
\end{definition}

\begin{remark}
$\prescript{L}{}{G}$ 的定义依赖于 $Q$ 的选取，尽管实践中用到它的地方往往与此无关。也可以尝试用 $\hat G\rtimes\Gal(\ol E/E)$，它不依赖任何选取；但那样一来，很难说清这是什么类型的对象（例如它不是代数群）。
\end{remark}

\begin{definition}[L-参数，第一种取法]
一个 L-参数是连续映射 $W_E\to\prescript{L}{}{G}(\CC)$，使下图交换：
\[
\begin{tikzcd}
W_E \arrow[rr] \arrow[dr] & & \prescript{L}{}{G}(\CC) \arrow[dl] \\
& Q &
\end{tikzcd}
\]
等价地，它是连续 1-上闭链 $W_E\to\hat G(\CC)$。
\end{definition}

\begin{remark}
连续性条件等价于要求该上闭链经过一个离散商 $W_E/I'$ 分解，其中 $I'\subset I_E$ 是开有限指数子群（因为 $\hat G(\CC)$ 的拓扑与 $I_E$ 的拓扑不相容）。

Deligne 观察到，更好的做法是同时记录一个单值算子（monodromy operator）。
\end{remark}

\begin{definition}[L-参数，第二种取法]
$\CC$ 上的一个 L-参数是一对 $(\phi,N)$，其中 $\phi\colon W_E\to\prescript{L}{}{G}(\CC)$ 如定义 26.3，而 $N\in\Lie\hat G\otimes\CC$，使得对所有 $w\in W_E$，
\[
\mathrm{Ad}(\phi(w))N=q^{|w|}N .\footnote{视正规化而定，也可能是 $q^{-|w|}N$。}
\]
\end{definition}

对 $G=\mathrm{GL}_n$，它们也称为 Weil--Deligne 表示。

\begin{definition}[L-参数，第三种取法]
$\CC$ 上的一个 L-参数是一对 $(\phi,r)$，其中 $\phi\colon W_E\to\prescript{L}{}{G}(\CC)$ 是连续群同态，$r\colon\mathrm{SL}_2\to\hat G/\CC$ 是代数表示，使得由
\[
\phi'(w):=\phi(w)\,r\begin{pmatrix}q^{|w|/2}&0\\0&q^{-|w|/2}\end{pmatrix},
\qquad
N:=(\Lie r)\begin{pmatrix}0&1\\0&0\end{pmatrix}
\]
定义的 $(\phi',N)$ 是定义 26.5 意义下的 L-参数。
\end{definition}

\begin{remark}
第一种、第二种与第三种取法各自自然地给出 L-参数的代数簇，它们彼此都不相同。

第二种与第三种取法下的参数在 $\hat G(\CC)$ 作用下双射，但概形结构不同。原因是在第二种取法中，$N\ne0$ 可以退化到 $N=0$。但由于 $\mathrm{SL}_2$ 的表示论半单，$\mathrm{SL}_2$ 的不同表示不能相互退化。我们希望有这种退化，所以第二种取法对我们而言是正确的。
\end{remark}

\subsection{Deligne 的动机}

固定一个同构 $\CC\cong\Ql$。

\begin{definition}[L-参数，取法二$'$]
$\Ql$ 上的一个 L-参数是连续群同态 $\phi_\ell\colon W_E\to\prescript{L}{}{G}(\Ql)$，它使下图交换：
\[
\begin{tikzcd}
W_E \arrow[rr, "\phi_\ell"] \arrow[dr] & & \prescript{L}{}{G}(\Ql) \arrow[dl] \\
& Q &
\end{tikzcd}
\]
或者换言之，它是连续 1-上闭链 $W_E\to\hat G(\Ql)$。
\end{definition}

取法二与取法二$'$在如下意义下等价。固定一个平凡化 $\Zl(1)\cong\Zl$ 与一个 Frobenius 元 $\Phi\in W_E$。这诱导一个收缩 $W_E\to I_E$；再与到 pro-$\ell$ 部分的投影以及该平凡化复合，得到 $t_\ell\colon W_E\to\Zl$。

\begin{theorem}[Grothendieck--Deligne]
沿用上述记号，任何使下图交换的 $\phi_\ell\colon W_E\to\prescript{L}{}{G}(\Ql)$
\[
\begin{tikzcd}
W_E \arrow[rr, "\phi_\ell"] \arrow[dr] & & \prescript{L}{}{G}(\Ql) \arrow[dl] \\
& Q &
\end{tikzcd}
\]
都形如 $\phi_\ell(w)=\phi(w)\exp(t_\ell(w)\cdot N)$，其中 $(\phi,N)$ 是定义 26.5 意义下唯一的 L-参数。
\end{theorem}

要点：同态 $W_E\to\mathrm{GL}_n(\Ql)$ 不必在开子群 $I'\subset I_E$ 上平凡；准确的说法是，可以找到这样的 $I'$，使它经过 $I'\twoheadrightarrow\Zl$ 分解。此时 $\Hom(\Zl,\mathrm{GL}_n(\Ql))$ 在开子群上由 $x\mapsto\exp(xN)$ 给出，其中 $N$ 是幂零矩阵。（论证是：特征值在开子群上必为 $1$，随后可以从对数中提取 $N$。）

Deligne 观察到，来自代数簇上同调的 Galois 表示中确实会出现这个单值算子 $N$。

\begin{warning}
定理 26.9 的双射确实依赖于辅助选取。我们将采用定义 26.8 作为 L-参数的定义。于是我们被迫在 $\Zl$ 上工作。
\end{warning}

\srcpage{139}

\subsection{L-参数的模空间}

我们的目标是构造 L-参数的模空间，即 $\Zl$ 上的概形 $Z^1(W_E,\hat G)$，使得它的 $A$-值点（对任意 $\Zl$-代数 $A$）是\srcpage{140}使下图交换的连续群同态
\[
\begin{tikzcd}
W_E \arrow[rr, "\phi"] \arrow[dr] & & \prescript{L}{}{G}(A) \arrow[dl] \\
& Q &
\end{tikzcd}
\]
即连续 1-上闭链 $W_E\to\hat G(A)$。（参考：[DHKM]。）

关于「连续性」，我们在 $A$ 上放什么拓扑？我们既不要离散拓扑，也不要 $\ell$-adic 拓扑，因为我们要把 $A=\Ql$ 包含进来。

\begin{definition}
任何 $\Zl$-模 $M$ 都可以赋予滤过余极限拓扑
\[
M=\varinjlim_{\Zl\text{-有限生成 }M'\subset M}\bigl(M',\ \ell\text{-adic 拓扑}\bigr).
\]
\end{definition}

\begin{remark}
用凝聚数学（condensed mathematics）的语言，相应的凝聚群可表为
\[
M=M^\delta\otimes_{\Zl^\delta}\Zl,
\]
或者用文字说：「把赋予离散拓扑的 $M$ 与赋予 $\ell$-adic 拓扑的 $\Zl$，在赋予离散拓扑的 $\Zl$ 上作张量」。这个等价性的内容是什么？由于所有构造都与滤过余极限相容，问题归结为看出：从投射有限集到这样的 $M$ 的任何映射，都经过某个有限生成的 $M'\subset M$ 分解。
\end{remark}

\begin{remark}
凝聚形式的一个好处是它能很好地容纳导出环（经典框架难以把握导出环与 $\ell$-adic 拓扑的混合）。然而可以证明，$Z^1(W_E,\hat G)$ 的导出版本已经是经典的。
\end{remark}

\begin{theorem}
存在 $\Zl$ 上的概形 $Z^1(W_E,\hat G)$ 参数化 $G$ 的 L-参数；它是 $\Zl$ 上有限型仿射概形的无交并，这些仿射概形平坦、是完全交，且维数为 $\dim G=\dim\hat G$。
\end{theorem}

\begin{remark}
对 $\hat G$ 的共轭作用取商，得到 Artin 叠「$\LocSys_{\hat G}$」。但它的连通分支有无穷多个，不同于几何 Langlands 中出现的经典 $\LocSys_{\hat G}$。
\end{remark}

\begin{proof}[证明概要]
第一点是：任何连续 1-上闭链 $\phi\colon W_E\to\hat G(A)$ 在野惯性子群的一个开子群 $P$ 上平凡。因此
\[
Z^1(W_E,\hat G)=\bigcup_P Z^1(W_E/P,\hat G),
\]
而且过渡映射是开闭浸入。

这是因为 $P$ 的元素必须映入 $\hat G(A)$ 中的 $p$ 次单位根，而后者是离散的。

因此只需证明：固定 $P$ 之后，所有 $Z^1(W_E/P,\hat G)$ 都是仿射的、平坦的、完全交，且维数为 $\dim\hat G$。

我们需要把（涉及连续性的）定义改写成更代数的东西。

为此有个技巧：我们如下构造一个稠密的离散子群 $W\subset W_E/P$。取 $\sigma\in W_E$ 为一个 Frobenius 元，$\tau\in I_E$ 为驯惯性的一个生成元。然后取由 $\sigma$、$\tau$ 以及野惯性子群生成的子群。这样得到短正合列
\[
1\to I\to W\to\ZZ\langle\sigma\rangle\to1
\]
与
\srcpage{141}
\[
1\to(\text{有限 }p\text{-群})\to I\to\ZZ[1/p]\langle\tau\rangle\to1 .
\]
要点是：$W$ 是 $W_E/P$ 的一个离散化。

断言：$Z^1(W_E/P,\hat G)\to Z^1(W,\hat G)$ 是同构。

左边显然是仿射代数簇，因为 $W_E/P$ 是有限生成离散群。我们可以把任何这样的表现变成 $Z^1(W_E/P,\hat G)$ 作为仿射代数簇的一个表现。

断言之证明：对任意 $A$，需要看出上闭链 $\phi_0\colon W\to\hat G(A)$ 有唯一的连续延拓 $\phi\colon W_E/P\to\hat G(A)$。唯一性由 $W$ 在 $W_E/P$ 中稠密是清楚的。需要证明存在性。为此可以把 $E$ 换成有限扩张，以消去那个有限 $p$-群。于是要看出的是：当表示
\[
\ZZ[1/p]\langle\tau\rangle\rtimes\ZZ\langle\sigma\rangle\to\mathrm{GL}_n(A)
\]
满足作用 $\sigma^{-1}\tau\sigma=\tau^q$ 时，映射 $\ZZ[1/p]\to\mathrm{GL}_n(A)$（把 $n$ 送到 $\im(\tau)^n$）连续延拓到 $\prod_{\ell\ne p}\Zl$。

注意 $\im(\tau)$ 与 $\im(\tau)^q$ 共轭，这蕴含所有特征值都是阶与 $p$ 互素的单位根。于是某个幂是幂幺的。而对幂幺矩阵，所有 $\Zl$-幂都有定义。

如上所述，该断言蕴含 $Z^1(W_E/P,\hat G)$ 是仿射概形。我们也可以用 Galois 上同调来度量形变理论。这清楚地表明它是完全交（因为 $W_E/P$ 的上同调维数为 $2$）。

为得到正确维数的平坦性，只需界住特殊纤维的维数。关键输入是 Lusztig 的一个定理：幂幺共轭类只有有限多个。要点是 $\tau$ 被送到一个幂幺矩阵。一旦固定共轭类，$\sigma$ 的像的选取就构成中心化子的一个主丛。这样便得到所需的维数界。
\end{proof}

\subsection{$\cO\bigl(Z^1(W_E/P,\hat G)\bigr)$ 的一个表现}

固定一个离散化 $W\subset W_E/P$。那么对任何映射 $F_n\to W$（$F_n$ 是自由群），我们得到映射
\[
Z^1(W_E/P,\hat G)=Z^1(W,\hat G)\to Z^1(F_n,\hat G)=\hat G^n .
\]

\begin{proposition}
有
\[
\varinjlim_{(n,F_n\to W)}\cO(\hat G^n)\xrightarrow{\ \sim\ }\cO\bigl(Z^1(W_E/P,\hat G)\bigr).
\]
\end{proposition}

\begin{remark}
余极限 $\varinjlim_{(n,F_n\to W)}\cO(\hat G^n)$ 是筛的（sifted），所以在模、代数等范畴中答案相同。它稍后会作为「巡游算子代数」出现。
\end{remark}

\begin{corollary}
映射
\[
\varinjlim_{(n,F_n\to W)}\cO(\hat G^n)^{\hat G}\longrightarrow
\cO\bigl(Z^1(W_E/P,\hat G)\bigr)^{\hat G}
=\cO\bigl(Z^1(W_E/P,\hat G)/\hat G\bigr)
\]
在谱上是泛同胚，且在反转 $\ell$ 之后是同构。
\end{corollary}

\begin{proof}
由于 $(-)^{\hat G}$ 在特征 $0$ 中正合，反转 $\ell$ 之后即可由上述命题推出。在正特征中，使用 Haboush 的「几何约化性」[Hab75]：它说明，在相差泛同胚的意义下，其行为如同半单情形。
\end{proof}

\begin{remark}
$\hat G$ 在 $\cO(\hat G^n)$ 上的作用通过等同 $\Hom(F_n,\hat G)\cong\hat G^n$ 实现，所以它是一种「扭共轭」，依赖于 $F_n$ 如何映入 $\hat G$。
\end{remark}

\srcpage{142}

\begin{definition}
我们称 $\cO\bigl(Z^1(W_E/P,\hat G)/\hat G\bigr)$ 为「谱 Bernstein 中心」。
\end{definition}

\begin{remark}
代数 $\varinjlim_{(n,F_n\to W)}\cO(\hat G^n)^{\hat G}$ 是「巡游算子代数」。
\end{remark}

\begin{theorem}
事实上，若 $Z(G)$ 连通且 $\ell$「不太小」（即 $\ell$ 是「好素数」），则上述映射是同构：
\begin{itemize}
\item 类型 $A$ 中的所有 $\ell$。
\item 类型 $A_n$、$B_n$、$C_n$、$D_n$、${}^{2}D_n$ 中所有 $\ell\ne2$。
\item 类型 ${}^{3}D_4$、${}^{6}D_4$、$E_6$、$E_7$、$F_4$、$G_2$、${}^{2}E_6$ 中所有 $\ell\ne2,3$。
\item 类型 $E_8$ 中所有 $\ell\ne2,3,5$。
\end{itemize}
这些条件蕴含 $\cO(Z^1)$ 有好滤过，从而更高的 $\hat G$-上同调消没，且不变量与任意基变换交换。
\end{theorem}

\begin{remark}
我们期望「好素数」假设是不必要的。$Z(G)$ 连通的假设不会给后续应用带来麻烦。
\end{remark}
